% ~7 min
\partframe{Part II}{Recursive Coalgebras and Hylomorphisms}

\begin{frame*}
  \frametitle{Why Not Just Write Quicksort?}
  \rocqblock{naive-qsort.v}
  \badblock{naive-qsort-err.txt}
  \speakernote{Rocq's structural check rejects it: the recursive calls are
    not on subterms of l. Program Fixpoint, Function or Equations accept it,
    but each function needs its own termination proof, and we get none of
    the laws.}
\end{frame*}

\begin{frame}
  \frametitle{Recursive Coalgebras (I)}
  \[ h = a \circ F\,h \circ c \qquad \text{with } F\,X = X,
     \text{ so } F\,h = h \]
  \[
    \renewcommand\arraystretch{1.6}
    \begin{array}{l@{\qquad}l@{\qquad}l}
      \uncover<2->{c = \op{id},\ a = \op{succ}} &
        \uncover<2->{h\,x = h\,x + 1} &
        \uncover<2->{\text{\textbf{no} solution}} \\
      \uncover<3->{c = a = \op{id}} &
        \uncover<3->{h\,x = h\,x} &
        \uncover<3->{\text{\textbf{every} $h$ is a solution}}
    \end{array}
  \]
  \begin{uncoverenv}<4->
    The culprit: $c = \op{id}$ never shrinks its input,
    $x \mapsto x \mapsto x \mapsto \cdots$

    \textbf{Recursive} $c$: every input reaches the leaves; then there is
    exactly one $h$, for every $a$.
  \end{uncoverenv}
  \slideref{J. Ad\'amek, S. Milius, L. S. Moss. \textbf{On Well-Founded and
    Recursive Coalgebras}. FoSSaCS 2020.\\
    V. Capretta, T. Uustalu, V. Vene. \textbf{Recursive Coalgebras from
    Comonads}. 2006.}
  \speakernote{Why hylo cannot be defined for arbitrary coalgebras. With
    the identity functor the hylo equation is just an equation on
    functions. succ: no number equals its successor, so no h. id: the
    equation says nothing, so every h fits. In both cases the call tree is
    the infinite chain x, x, x, ... Recursive coalgebras are those whose
    unfolds are finite trees (they terminate on every input); for them the
    solution exists and is unique.}
\end{frame}

\begin{frame*}
  \frametitle{Recursive Coalgebras (and II)}
  \rocqblock{recf.v}
  \begin{center}
    \begin{tikzpicture}[alt={A proof of RecF c x is a finite tree that
        mirrors the call tree: one node for x, with a sub-proof for each
        child of c x, down to leaves with no valid positions.},
      level distance = 9mm, sibling distance = 22mm,
      every node/.style = {box, font = \small}]
      \node {$\op{RecF}\,c\,x$}
        child {node {$\op{RecF}\,c\,x_1$}
          child {node {leaf}}}
        child {node {$\op{RecF}\,c\,x_2$}
          child {node {leaf}} child {node {leaf}}};
    \end{tikzpicture}
  \end{center}
  Equivalent to \code{well\_founded (child c)}, but the induction
  hypothesis is \emph{per position}, as in \code{fmap}.
  \speakernote{RecF: c terminates on x if it terminates on every child of
    c x. A proof is a finite tree shaped like the call tree. In practice it
    comes from a well-founded measure that c decreases (mk\_wf\_coalg),
    e.g. list length for quicksort. Why not standard accessibility? RecF c
    x is equivalent to Acc (child c) x, where child y x says y is at some
    position of c x (RecP\_well\_founded, so Equations and the standard
    library can use our proofs). Same hylo, same extracted code, same
    termination proofs. The difference is in correctness proofs: RecF's
    induction hypothesis is forall e, P (cont (c x) e), exactly what the
    unfolded hylo exposes, so it is applied at a position directly; with
    Acc it is forall y, child y x -> P y, and each use must build the
    witness and unfold child.}
\end{frame*}

\begin{frame*}
  \frametitle{Recursive Hylomorphisms}
  Recursion on the termination proof (simplified):
  \rocqblock[width = .8\textwidth]{hylo-def.v}
  \callout<2->{4,2.4}{termination proof\\ (in \code{Prop}: erased)}{recf}
  \callout<3->{13,2.4}{structurally smaller proof}{He}
  \speakernote{Structural recursion on the proof H, not on x. The proof is
    in Prop, so extraction erases it. RecF\_inv gives the structurally
    smaller proof for each child. The real definition needs a dependent
    match; this is simplified. The result does not depend on the proof.}
\end{frame*}

\begin{frame*}
  \frametitle{Uniqueness and the Laws}
  \rocqblock{hylo-laws.v}
  \popup<2>(10,60)[.9][greenbox]{%
    Proofs follow the paper proofs:
    \[
      \begin{tikzcd}[alt={Proof of left fusion: the squares for the hylo of
          a and c and for the homomorphism h paste into a rectangle whose
          top is h after the hylo, so by uniqueness it equals the hylo of b
          and c.}, column sep = 24mm, row sep = 8mm]
        F\,C \arrow[r, "F\,\hylo{a}{c}"] \& F\,A \arrow[r, "F\,h"] \&
          F\,B \\
        C \arrow[u, "c"] \arrow[r, "\hylo{a}{c}"'] \& A \arrow[u, "a"',
          leftarrow] \arrow[r, "h"'] \& B \arrow[u, "b"', leftarrow]
      \end{tikzcd}
    \]}
  \speakernote{hylo\_uniq is the universal property; the three laws
    follow from it. The Rocq proof of fusion is this diagram: paste the
    squares, apply uniqueness; four rewrite steps.}
\end{frame*}

\begin{frame*}
  \frametitle{Folds after Unfolds}
  \rocqblock{cata-ana.v}
  \[
    \begin{tikzcd}[alt={The specification goes from A to LFix F by rana c,
        then to B by cata a. The hylomorphism goes directly from A to B and
        never builds the tree.}, column sep = 30mm]
      A \arrow[r, "\code{rana c}"] \arrow[rr, bend right = 25,
        "\code{hylo a c}"', dashed] \& \code{LFix F}
        \arrow[r, "\code{cata a}"] \& B
    \end{tikzcd}
  \]
  \speakernote{rana builds the call tree (a hylo into the initial
    algebra). Specifications are written as fold after unfold; the
    calculated program never builds the tree. Correctness proofs: induction
    on termination, unroll the hylo once; they carry over along every
    calculation step.}
\end{frame*}

\begin{frame*}
  \frametitle{No Axioms}
  \outblock{assumptions.txt}
  \speakernote{Print Assumptions output. The core laws depend on nothing:
    no functional extensionality, no axiom K, no proof irrelevance. The
    examples with machine integers only depend on Rocq's own primitive
    integer axioms, which extract to native OCaml integers.}
\end{frame*}
